\documentclass[12pt]{ximera}
\title{Homework 9: Hamilton's and Jefferson's Methods}
\begin{document}
\begin{abstract}
Sections 4.1--4.3: standard divisors and quotas, and Hamilton's and Jefferson's methods
applied to hospital nurses and to seventeen residence halls --- where Jefferson's method
breaks the quota rule.
\end{abstract}
\maketitle

\section*{Sections 4.1--4.3: Hamilton's and Jefferson's Methods}

\begin{problem}
A hospital recently hired $25$ nurses and wants to assign them proportionally across
four departments based on average monthly patient load:
\begin{center}
\begin{tabular}{rl}
Pediatrics: & $328$ patients\\
Surgery: & $177$ patients\\
Maternity: & $132$ patients\\
Emergency: & $108$ patients
\end{tabular}
\end{center}

\begin{prompt}
\textbf{(a)} Calculate the standard divisor.

$\answer[tolerance=0.005]{29.8}$

What does it represent here?
\begin{multipleChoice}
\choice[correct]{The number of patients per nurse}
\choice{The number of nurses per department}
\choice{The number of patients per department}
\end{multipleChoice}

\begin{feedback}
The total load is $328+177+132+108=745$ patients, so
$SD=\frac{745}{25}=29.8$ --- about $29.8$ patients per nurse.
\end{feedback}
\end{prompt}

\begin{prompt}
\textbf{(b)} Compute the standard quota for each department, to two decimal places.

Pediatrics $\answer[tolerance=0.005]{11.01}$ \quad Surgery $\answer[tolerance=0.005]{5.94}$
\quad Maternity $\answer[tolerance=0.005]{4.43}$ \quad Emergency $\answer[tolerance=0.005]{3.62}$

\begin{feedback}
Divide each load by $29.8$: $\frac{328}{29.8}=11.01$, $\frac{177}{29.8}=5.94$,
$\frac{132}{29.8}=4.43$, $\frac{108}{29.8}=3.62$.
\end{feedback}
\end{prompt}

\begin{prompt}
\textbf{(c)} How many nurses does each department receive under Hamilton's method?

Pediatrics $\answer{11}$ \quad Surgery $\answer{6}$ \quad Maternity $\answer{4}$ \quad
Emergency $\answer{4}$

\begin{feedback}
The lower quotas $11, 5, 4, 3$ total $23$, leaving $2$ nurses. The largest fractional
parts are Surgery ($.94$) and Emergency ($.62$), so each gets one more:
$11, 6, 4, 4$.
\end{feedback}
\end{prompt}
\end{problem}

\begin{problem}
Suppose someone on the hiring team argues that they should use Jefferson's method to
apportion the nurses instead.

\begin{prompt}
\textbf{(a)} Which of these is an appropriate modified divisor?

\begin{multipleChoice}
\choice{$29.8$}
\choice{$28.0$}
\choice[correct]{$27.2$}
\choice{$26.5$}
\end{multipleChoice}

\begin{hint}
Rounding the standard quotas down gives only $23$ nurses, so the divisor has to shrink.
\end{hint}

\begin{feedback}
Any divisor greater than $27$ and at most about $27.33$ works. At $28.0$ the rounded
quotas total only $24$; at $26.5$ they total $26$. With $27.2$ they total exactly $25$.
\end{feedback}
\end{prompt}

\begin{prompt}
\textbf{(b)} Using $d=27.2$, compute each modified quota, to two decimal places.

Pediatrics $\answer[tolerance=0.005]{12.06}$ \quad Surgery $\answer[tolerance=0.005]{6.51}$
\quad Maternity $\answer[tolerance=0.005]{4.85}$ \quad Emergency $\answer[tolerance=0.005]{3.97}$

\begin{feedback}
$\frac{328}{27.2}=12.06$, $\frac{177}{27.2}=6.51$, $\frac{132}{27.2}=4.85$,
$\frac{108}{27.2}=3.97$.
\end{feedback}
\end{prompt}

\begin{prompt}
\textbf{(c)} How many nurses does each department receive under Jefferson's method?

Pediatrics $\answer{12}$ \quad Surgery $\answer{6}$ \quad Maternity $\answer{4}$ \quad
Emergency $\answer{3}$

\begin{feedback}
Rounding down: $12+6+4+3=25$. \checkmark Compared with Hamilton's method, one nurse
moves from Emergency, the smallest department, to Pediatrics, the largest --- the usual
tilt of Jefferson's method. (Pediatrics's upper quota is $12$, so no quota rule is
broken here.)
\end{feedback}
\end{prompt}
\end{problem}

\begin{problem}
The residence halls at a university have the following population capacities,
totaling $6558$ students:

\begin{center}
\begin{tabular}{r|c|c|c|c|c|c}
\textbf{Hall} & Aspen & Cedar & Pine & Birch & Willow & Larch \\\hline
\textbf{Pop.} & 270 & 270 & 331 & 184 & 544 & 544 \\
\end{tabular}

\smallskip
\begin{tabular}{r|c|c|c|c|c|c}
\textbf{Hall} & Spruce & Poplar & Hickory & Maple & Linden & Elm \\\hline
\textbf{Pop.} & 99 & 136 & 107 & 496 & 321 & 202 \\
\end{tabular}

\smallskip
\begin{tabular}{r|c|c|c|c|c}
\textbf{Hall} & Walnut & Oak & Alder & Chestnut & Sycamore \\\hline
\textbf{Pop.} & 770 & 235 & 116 & 713 & 1220 \\
\end{tabular}
\end{center}

The university has hired $220$ RAs altogether. Consider this situation using both
Hamilton's and Jefferson's methods.

\begin{prompt}
\textbf{(a)} Compute the standard divisor, to two decimal places.

$\answer[tolerance=0.005]{29.81}$

\smallskip
What does that value represent in this situation?
\begin{multipleChoice}
\choice[correct]{The number of students per RA}
\choice{The number of RAs per residence hall}
\choice{The number of residence halls per RA}
\end{multipleChoice}

\begin{feedback}
The standard divisor is the total population divided by the number of seats:
\[ SD = \frac{6558}{220} \approx 29.81. \]
It represents the number of \emph{students per RA} --- roughly one RA for every $30$
students.
\end{feedback}
\end{prompt}

\begin{prompt}
\textbf{(b)} Compute the standard quota of each of these residence halls, to two
decimal places. (A hall's quota is its population divided by the standard divisor.)

Sycamore: $\answer[tolerance=0.005]{40.93}$ \qquad
Walnut: $\answer[tolerance=0.005]{25.83}$ \qquad
Chestnut: $\answer[tolerance=0.005]{23.92}$

\smallskip
Willow: $\answer[tolerance=0.005]{18.25}$ \qquad
Maple: $\answer[tolerance=0.005]{16.64}$ \qquad
Spruce: $\answer[tolerance=0.005]{3.32}$

\begin{feedback}
Dividing each population by $29.81$:
\[
\begin{array}{llll}
\text{Aspen } 9.06 & \text{Cedar } 9.06 & \text{Pine } 11.10 & \text{Birch } 6.17\\
\text{Willow } 18.25 & \text{Larch } 18.25 & \text{Spruce } 3.32 & \text{Poplar } 4.56\\
\text{Hickory } 3.59 & \text{Maple } 16.64 & \text{Linden } 10.77 & \text{Elm } 6.78\\
\text{Walnut } 25.83 & \text{Oak } 7.88 & \text{Alder } 3.89 & \text{Chestnut } 23.92\\
\text{Sycamore } 40.93 & & &
\end{array}
\]
These sum to $220$, as the standard quotas always must.
\end{feedback}
\end{prompt}

\begin{prompt}
\textbf{(c)} What is the total of the (standard) \emph{lower} quotas?

$\answer{211}$

\smallskip
How many surplus RAs remain to be handed out? $\answer{9}$

\begin{hint}
Round every quota \emph{down} and add. The shortfall from $220$ is the surplus.
\end{hint}

\begin{feedback}
Rounding every quota down and adding gives $211$. Since $220$ RAs must be placed, there
is a surplus of
\[ 220 - 211 = 9 \]
RAs still to assign.
\end{feedback}
\end{prompt}

\begin{prompt}
\textbf{(d)} Under Hamilton's method, the surplus goes to the halls with the largest
fractional parts. Which halls receive an extra RA?

\begin{selectAll}
\choice{Aspen}
\choice{Cedar}
\choice{Pine}
\choice{Birch}
\choice{Willow}
\choice{Larch}
\choice{Spruce}
\choice{Poplar}
\choice[correct]{Hickory}
\choice[correct]{Maple}
\choice[correct]{Linden}
\choice[correct]{Elm}
\choice[correct]{Walnut}
\choice[correct]{Oak}
\choice[correct]{Alder}
\choice[correct]{Chestnut}
\choice[correct]{Sycamore}
\end{selectAll}

\begin{feedback}
Ranking the halls by fractional part, largest first:
\[
\begin{array}{lll}
\text{Sycamore } .927 & \text{Chestnut } .919 & \text{Alder } .891\\
\text{Oak } .884 & \text{Walnut } .831 & \text{Elm } .776\\
\text{Linden } .769 & \text{Maple } .639 & \text{Hickory } .590\\
\end{array}
\]
Those are the top nine, so each of them gets one extra RA. The next in line is Poplar
at $.562$, which just misses.

The final Hamilton apportionment is:
\[
\begin{array}{llll}
\text{Aspen } 9 & \text{Cedar } 9 & \text{Pine } 11 & \text{Birch } 6\\
\text{Willow } 18 & \text{Larch } 18 & \text{Spruce } 3 & \text{Poplar } 4\\
\text{Hickory } 4 & \text{Maple } 17 & \text{Linden } 11 & \text{Elm } 7\\
\text{Walnut } 26 & \text{Oak } 8 & \text{Alder } 4 & \text{Chestnut } 24\\
\text{Sycamore } 41 & & &
\end{array}
\]
which totals $220$. \checkmark

Notice Hickory (population $107$) receives $4$ RAs while Poplar (population $136$)
receives only $4$ as well --- and Alder, with just $116$ students, also gets $4$. Small
halls can do disproportionately well or badly on the strength of a fractional part.
\end{feedback}
\end{prompt}
\begin{prompt}
\textbf{(e)} Now use Jefferson's method (round every modified quota \emph{down}). Which
modified divisor works?

\begin{multipleChoice}
\choice{$29.81$ (the standard divisor)}
\choice{$29.2$}
\choice[correct]{$28.7$}
\choice{$28.2$}
\end{multipleChoice}

\begin{hint}
Rounding the standard quotas down gives only $211$ RAs, so the divisor must shrink until
exactly $220$ survive the rounding. Try each option.
\end{hint}

\begin{feedback}
Any divisor from about $28.64$ to $28.86$ works. With $d=28.7$, rounding each
$\frac{\text{pop}}{28.7}$ down gives a total of exactly $220$. At $29.2$ the total is
still short, and at $28.2$ it overshoots.
\end{feedback}
\end{prompt}

\begin{prompt}
\textbf{(f)} Under Jefferson's method, which two halls come out differently from
Hamilton's method?

Sycamore receives $\answer{42}$ RAs (Hamilton gave $41$), and Hickory receives
$\answer{3}$ (Hamilton gave $4$). Every other hall is the same.

\begin{feedback}
With $d=28.7$: Sycamore gets $\lfloor 1220/28.7\rfloor=\lfloor 42.51\rfloor=42$, and
Hickory gets $\lfloor 107/28.7\rfloor=\lfloor 3.73\rfloor=3$.

Sycamore's standard quota was $40.93$, so its upper quota is $41$ --- and Jefferson's
method gives it $42$. This is a \textbf{quota rule violation}, in the direction
Jefferson's method always leans: toward the largest hall, at the expense of a small one.
\end{feedback}
\end{prompt}
\end{problem}

\end{document}
